\section{Research Questions and Contributions}

We organize the evaluation around seven research questions.

\paragraph{RQ1 --- Accuracy and robustness.}
How do local ASR engines and checkpoints compare when the same product pipeline is evaluated on
both clean and harder speech?

\paragraph{RQ2 --- Reproducibility.}
When identical audio is decoded repeatedly, are the resulting hypotheses stable? If not, which
error components change?

\paragraph{RQ3 --- Decoder mechanism.}
How do beam width, temperature fallback, and previous-text conditioning interact with checkpoint
size and catastrophic continuation behavior?

\paragraph{RQ4 --- Hardware and platform portability.}
How much do decoded outputs and installation constraints vary across processor architecture and
operating system?

\paragraph{RQ5 --- Interactive latency assumptions.}
Do synthetic onset padding and streaming partial decoding provide their intended benefits under a
CPU-only interactive workload?

\paragraph{RQ6 --- Meeting diarization.}
Does the system's diarization configuration transfer to real annotated meetings, and how do
clustering threshold and a speaker-count ceiling affect different aspects of performance?

\paragraph{RQ7 --- Reproducibility infrastructure.}
Does retaining commands, corpus identity, provenance, failed rows, and displaced benchmark
results materially change the conclusions reached from an experimental campaign?

\subsection{Measured contributions}

This work makes the following measured systems contributions.

\begin{enumerate}
  \item A common-harness CPU comparison of multiple local speech-recognition engines and
  checkpoints through the same shipping \YazSes{} engine interface, evaluated on both clean and
  harder LibriSpeech conditions.

  \item A characterization of a repeated-decoding failure observed with \texttt{large-v3}. In
  the measured hard-speech workload, variation in aggregate \wer{} is localized to insertions and
  runaway continuation while substitutions, deletions, and correct-word hits remain unchanged
  across the analyzed repeats.

  \item Evidence that decoder configuration is checkpoint-dependent. In particular, the measured
  effect of previous-text conditioning changes across the Whisper checkpoint ladder rather than
  following a single monotonic rule.

  \item Decoding results across Linux x86-64, Linux arm64, macOS arm64, and Windows x86-64 using
  a common subset, together with a separate dependency-resolution analysis.

  \item A full AMI-test-split Meeting Mode evaluation showing that the previous clustering
  operating point is severely mismatched to long real meetings, and that speaker-count
  correctness and diarization error are distinct objectives.

  \item Negative results with direct product implications: synthetic pre-speech silence does not
  recover clipped speech, generic CPU streaming does not consistently reduce final latency, and
  simple centroid-based post-hoc cluster merging does not provide a safe repair strategy on the
  measured AMI data.

  \item An auditable benchmark archive that ties results to command-line arguments, corpus
  identity, machine/library provenance, and retained displaced reruns.
\end{enumerate}

We deliberately distinguish these measured contributions from claims of literature priority.
Priority language remains deferred until the current related-work review in repository
issue \#507 is complete.